\chapter{Commutateurs et dispositifs de couche 1}\label{ch:nw:switches-and-layer-1-devices}

Un commutateur store-and-forward ne peut pas réexpédier une trame de 1\,518~octets avant d'en avoir reçu la totalité~: à
\qty{10}{\giga\bit\per\second}, cela représente \qty{1.2}{\micro\second} d'attente, quoi que fasse ensuite le commutateur. Un \omterm{def:nw:switches-and-layer-1-devices:l1}{commutateur de couche 1}
vendu aux sociétés de trading fait transiter les données entre ses ports en \qty{4}{\nano\second}, avec moins de 100~picosecondes de gigue, et, selon les
mots de sa fiche technique, il «~ne met pas en mémoire tampon ni en file d'attente les données~»~: il ne regarde jamais la trame. Entre les deux se trouvent les
commutateurs cut-through, qui ne lisent que les premiers octets. Choisir parmi eux, port par port, constitue l'essentiel de ce que
signifie concevoir le réseau à l'intérieur d'une cage de colocation.

Ce chapitre suit une trame à travers chaque type d'équipement, puis à travers les équipements passifs qui permettent à une société de surveiller son propre
trafic, et se termine par trois conceptions d'une même petite cage~: la latence de chaque chemin, les équipements que traverse chaque chemin, et les
seuls équipements dont la défaillance isolerait un serveur.

\section{Commutation store-and-forward et cut-through}

\begin{definition}[Commutation store-and-forward et commutation cut-through]\label{def:nw:switches-and-layer-1-devices:modes}
Un commutateur qui utilise la \emph{commutation store-and-forward}\index{commutation store-and-forward} reçoit une trame en entier, vérifie sa
séquence de contrôle de trame, et ne commence qu'alors à l'émettre sur le port de sortie. Un commutateur qui utilise la \emph{commutation
cut-through}\index{commutation cut-through} commence à émettre une trame dès qu'il a lu les champs d'en-tête dont il a besoin pour choisir
le port de sortie, alors que le reste de la trame arrive encore.
\end{definition}

\begin{proposition}[Ce qu'ajoute un commutateur]\label{prop:nw:switches-and-layer-1-devices:add}
Soit une trame de $\ell_{\mathrm f}$ octets qui arrive à $R_{\mathrm L}$ sur un port dont la sortie est libre et de même débit, et soit
$\tau$ le temps de décision et de transfert interne du commutateur. Un commutateur store-and-forward ajoute
$8\ell_{\mathrm f}/R_{\mathrm L} + \tau$ à la latence de la trame~; un commutateur cut-through qui lit les $h$ premiers octets ajoute
$8h/R_{\mathrm L} + \tau$, indépendamment de la longueur de la trame. Si le port de sortie est occupé, ou s'il est plus rapide que l'entrée (il manquerait alors
de bits au milieu de la trame), un commutateur cut-through doit mettre la trame en mémoire tampon et revient à la \omterm{def:nw:switches-and-layer-1-devices:modes}{commutation store-and-forward} pour
celle-ci.
\end{proposition}

\begin{proof}
Le dernier bit quitte l'émetteur à $8\ell_{\mathrm f}/R_{\mathrm L}$ et atteint l'extrémité distante une sérialisation après que le
premier bit a quitté le dernier équipement. Un commutateur store-and-forward émet son premier bit $\tau$ après l'arrivée de son dernier bit, de sorte que la sérialisation
est payée deux fois, une fois à l'entrée du commutateur et une fois à sa sortie~; un commutateur cut-through émet son premier bit $\tau$ après l'arrivée du $h$-ième octet,
et les deux sérialisations se recouvrent. Si la sortie est occupée, les bits doivent attendre~; si elle est plus rapide que l'entrée, elle émettrait
plus vite que les bits n'arrivent, et une trame ne peut pas être suspendue au milieu, de sorte qu'elle doit d'abord être stockée.
\end{proof}

\begin{omfigure}
\begin{tikzpicture}[font=\footnotesize,fr/.style={draw,minimum height=0.45cm,anchor=west,inner sep=0pt}]
\node[anchor=east] at (-0.2,2.1) {entrée (les deux)};
\node[fr,fill=omDef!20,minimum width=5.0cm] at (0,2.1) {trame qui arrive};
\node[anchor=east] at (-0.2,1.2) {sortie, cut-through};
\node[fr,fill=omMeth!25,minimum width=5.0cm] at (1.3,1.2) {trame qui part};
\node[anchor=east] at (-0.2,0.3) {sortie, store-and-forward};
\node[fr,fill=omThm!20,minimum width=5.0cm] at (5.8,0.3) {trame qui part};
\draw[densely dashed,gray] (0.5,2.45) -- (0.5,-0.3);
\draw[densely dashed,gray] (5.0,2.45) -- (5.0,-0.3);
\draw[<->] (0.5,-0.25) -- node[below,font=\scriptsize]{lecture de l'en-tête, puis $\tau$} (1.3,-0.25);
\draw[<->] (5.0,-0.25) -- node[below,font=\scriptsize]{$\tau$} (5.8,-0.25);
\draw[->] (0,-0.9) -- (11.4,-0.9) node[right]{temps};
\end{tikzpicture}

\omcaption{La même trame à travers un commutateur cut-through puis un commutateur store-and-forward. Les deux attendent le temps interne
$\tau$~du commutateur~; le commutateur cut-through commence après l'en-tête (pointillé de gauche), le commutateur store-and-forward après le dernier bit
(pointillé de droite), de sorte qu'il ajoute une sérialisation complète de la trame.}\label{fig:nw:switches-and-layer-1-devices:modes}
\end{omfigure}

Un commutateur store-and-forward élimine une trame dont la séquence de contrôle est erronée avant de la faire progresser~; un commutateur cut-through a déjà
envoyé l'essentiel de la trame lorsqu'il s'en aperçoit, et ne peut que marquer la fin de la trame comme erronée. C'est tout le prix du cut-through, et
dans une cage de colocation, où les liaisons sont courtes et les erreurs rares, c'est un prix que les sociétés acceptent de payer.

\begin{example}[Un ordre, trois commutateurs]\label{ex:nw:switches-and-layer-1-devices:order}
Une trame d'ordre de 100~octets à \qty{10}{\giga\bit\per\second}. Un commutateur store-and-forward avec un temps interne de \qty{500}{\nano\second}
(une valeur de modèle, et non celle d'une fiche technique) ajoute $80 + 500 = \qty{580}{\nano\second}$~; un commutateur cut-through qui lit 64~octets ajoute
\qty{51.2}{\nano\second} plus son temps interne~; un équipement de couche 1 n'ajoute que ses quelques nanosecondes propres. Pour une trame de 1\,518~octets,
le coût du commutateur store-and-forward passe à \qty{1\,714}{\nano\second}, celui des autres ne change pas
(\cref{fig:nw:switches-and-layer-1-devices:curve}).
\end{example}

\section{Équipements de couche 1~: duplication, multiplexage et brassage}

\begin{definition}[Commutateur de couche 1, duplication, multiplexeur]\label{def:nw:switches-and-layer-1-devices:l1}
Un \emph{commutateur de couche 1}\index{commutateur de couche 1} relie des ports au niveau du signal électrique ou optique~: il régénère
les bits d'une entrée sur les sorties qui lui sont configurées, sans décoder les trames, de sorte qu'il ne lit ni adresses ni ne met rien en file d'attente. Une
\emph{duplication de couche 1}\index{duplication de couche 1} (fan-out) est une configuration qui recopie une entrée vers plusieurs sorties à la fois. Un
\emph{multiplexeur de couche 1}\index{multiplexeur de couche 1} fusionne les trames de plusieurs entrées sur une seule sortie~; contrairement à une duplication, il doit
reconnaître les limites des trames et retenir une trame qui arrive pendant qu'une autre est émise.
\end{definition}

La duplication est l'équipement naturel pour les données de marché~: le flux de la plateforme arrive une seule fois sur l'interconnexion et doit atteindre chaque
serveur qui y négocie, tous au même instant, sans rien avoir à décider. Le multiplexeur est l'équipement naturel pour les ordres~: de nombreux
serveurs, une seule remise d'ordres, et aucune décision de routage au-delà de «~tout vers ce port~». Un multiplexeur ne peut pas être aussi rapide qu'une duplication,
parce que deux serveurs peuvent émettre en même temps.

\begin{proposition}[Contention au multiplexeur]\label{prop:nw:switches-and-layer-1-devices:mux}
Deux trames de temps de sérialisation $s$ atteignent la sortie d'un multiplexeur à des instants qui diffèrent de $D$. Si $|D| < s$, la trame la plus tardive attend
$s - |D|$~; sinon, aucune n'attend. Si la réponse de chaque serveur à un événement commun arrive avec un délai normal indépendant,
d'écart type $\sigma$, $D$ est normal d'écart type $\sigma\sqrt2$, et le plus tardif de deux ordres attend avec
la probabilité $\P(|D| < s) = 2\Phi\bigl(s/(\sigma\sqrt2)\bigr) - 1$.
\end{proposition}

\begin{proof}
La sortie n'émet qu'une trame à la fois~: la trame qui arrive en second commence lorsque la première est terminée, $s$ après le
début de la première, ou à sa propre arrivée si celle-ci est plus tardive. La différence de deux variables normales indépendantes de variance $\sigma^2$ est
normale de variance $2\sigma^2$.
\end{proof}

\begin{example}[Deux serveurs, une remise]\label{ex:nw:switches-and-layer-1-devices:two}
Un ordre de 100~octets occupe \qty{96}{\nano\second} de temps de ligne à \qty{10}{\giga\bit\per\second}. Si deux serveurs de la société
répondent au même événement de marché avec des délais d'écart type \qty{50}{\nano\second}, le second ordre attend avec la probabilité
$2\Phi(96/70.7) - 1 = 0.83$. Les serveurs de la société se font concurrence au multiplexeur au moment qui compte
le plus, et celui dont la réponse a été plus rapide de quelques nanosecondes passe en premier.
\end{example}

Le tutoriel du chapitre étudie cette course pour deux, quatre et huit serveurs (\cref{fig:nw:switches-and-layer-1-devices:mux}). Avec quatre
serveurs répondant à moins de \qty{50}{\nano\second} les uns des autres, un ordre attend en moyenne \qty{93}{\nano\second}, et le dernier
\qty{186}{\nano\second}, soit plus que les \qty{39}{\nano\second} du multiplexeur lui-même~; si l'on étale les réponses sur une microseconde, la moyenne
tombe à \qty{6}{\nano\second}. Plus les serveurs d'une société sont synchronisés, plus ils se font attendre les uns les autres.

\begin{dated}{2026-09}{Latences relevées dans les fiches techniques}\label{dat:nw:switches-and-layer-1-devices:devices}
Les \omterm{def:nw:switches-and-layer-1-devices:l1}{commutateurs de couche 1} 7130 Connect Series d'Arista relaient les trames entre ports en \qty{4}{\nano\second} (\qty{6}{\nano\second} pour le
modèle à 96~ports), avec une gigue (jitter) inférieure à \qty{100}{\pico\second} et la même latence pour la recopie d'un port vers plusieurs. Le MetaMux d'Arista
multiplexe plusieurs entrées sur une sortie en \qty{39}{\nano\second} en moyenne (\qty{42}{\nano\second} sur un autre
modèle matériel)~; «~lorsque plusieurs paquets arrivent au multiplexeur au même instant~», il les met en file dans des tampons d'entrée, et sans
contention, sa latence varie de $\pm\qty{7}{\nano\second}$. Le commutateur cut-through Nexus 3548-X de Cisco annonce des latences «~aussi basses que~»
\qty{250}{\nano\second}, \qty{200}{\nano\second} en mode warp, et recopie un port vers un nombre quelconque d'autres ports dans son mode warp SPAN,
le commutateur annonce «~aussi basses que~» \qty{50}{\nano\second}.
\end{dated}

\begin{omfigure}
\begin{tikzpicture}
\begin{axis}[width=11.5cm,height=5.4cm,ymode=log,xmin=0,xmax=1600,ymin=2,ymax=4000,
  xlabel={taille de la trame (octets)},ylabel={latence ajoutée (ns, échelle log)},
  ytick={3,10,30,100,300,1000,3000},yticklabels={3,10,30,100,300,1\,000,3\,000},
  xtick={64,256,512,1024,1518},xticklabels={64,256,512,1\,024,1\,518},
  tick label style={font=\small},label style={font=\small},grid=major,grid style={gray!25},
  legend columns=2,legend style={font=\footnotesize,at={(0.5,-0.3)},anchor=north,draw=none,fill=none,
  /tikz/every even column/.append style={column sep=8pt}},table/col sep=comma]
\addplot[omThm,very thick,mark=*,mark size=1.4pt] table[x=frame,y=store_forward]{figdata/networks/02-switches-and-layer-1-devices/forwarding.csv};
\addplot[omMeth,very thick,mark=square*,mark size=1.4pt] table[x=frame,y=cut_through]{figdata/networks/02-switches-and-layer-1-devices/forwarding.csv};
\addplot[omExo,very thick,mark=triangle*,mark size=1.8pt] table[x=frame,y=mux]{figdata/networks/02-switches-and-layer-1-devices/forwarding.csv};
\addplot[omDef,very thick,mark=diamond*,mark size=1.8pt] table[x=frame,y=l1]{figdata/networks/02-switches-and-layer-1-devices/forwarding.csv};
\legend{commutation store-and-forward (modèle),commutation cut-through,multiplexeur,commutateur de couche 1}
\end{axis}
\end{tikzpicture}

\omcaption{Latence ajoutée par chaque appareil à une trame à \qty{10}{\giga\bit\per\second}, sortie inactive, selon la taille de la trame. Seul le
commutateur store-and-forward croît avec la trame (son temps interne de \qty{500}{\nano\second} est une valeur de modèle)~; les autres utilisent les
valeurs des fiches techniques de l'\cref{dat:nw:switches-and-layer-1-devices:devices}. Données~: \texttt{fig\_cage.py}.}\label{fig:nw:switches-and-layer-1-devices:curve}
\end{omfigure}

Un \omterm{def:nw:switches-and-layer-1-devices:l1}{commutateur de couche 1} est aussi un panneau de brassage piloté par logiciel~: n'importe quel port peut être relié à n'importe quel autre, ou à plusieurs, par
un changement de configuration plutôt qu'un technicien qui déplace un câble. Les sociétés s'en servent pour faire passer un serveur de la ligne d'une plateforme à celle d'une autre,
pour remplacer un appareil défaillant par un appareil de réserve, ou pour copier une liaison vers un dispositif de capture pendant une heure.

\section{Taps réseau et recopies de ports}

Une société qui veut savoir ce qu'elle a envoyé et reçu, et quand, doit copier son trafic sans le perturber.

\begin{definition}[Tap réseau, recopie de port]\label{def:nw:switches-and-layer-1-devices:tap}
Un \emph{tap réseau}\index{tap réseau} est un appareil inséré dans une liaison qui copie son trafic vers des ports de supervision~; un tap optique passif
le fait en divisant la lumière de chaque fibre, en envoyant une part fixe (son taux de division) vers le moniteur, et n'a besoin
d'aucune alimentation et n'ajoute que la longueur de sa fibre. La \emph{recopie de port}\index{recopie de port} est la copie, par un commutateur, du trafic de certains de
ses ports vers un autre port, par logiciel ou dans son matériel de commutation.
\end{definition}

Un tap se place sur la liaison elle-même~: il ne peut ni perdre, ni réordonner, ni retarder la copie par rapport à l'original, et il continue de fonctionner quand le
commutateur ne fonctionne plus. Son coût tient à la lumière qu'il prélève (une répartition 50/50 divise par deux la puissance qui atteint l'autre extrémité, ce qu'une longue interconnexion (cross-connect)
peut ne pas pouvoir se permettre) et un port de supervision par sens. Un port de recopie est gratuit, puisque le commutateur est déjà là, mais
c'est un port de commutateur comme les autres~: recopier les deux sens d'une liaison très chargée à \qty{10}{\giga\bit\per\second} envoie jusqu'à
\qty{20}{\giga\bit\per\second} vers un port à \qty{10}{\giga\bit\per\second}, qui en perd la moitié quand les deux sens sont saturés, précisément dans les
rafales qu'on voulait observer. Les dispositifs de capture (chapitre~5) se placent donc sur des taps et sur les sorties de recopie de couche 1, qui copient au
\omterm{def:nw:networking-for-trading:ser}{débit de ligne}.

\section{Concevoir le réseau à l'intérieur d'une cage}

\begin{definition}[Réseau d'administration hors bande, câble à connexion directe]\label{def:nw:switches-and-layer-1-devices:oob}
Un \emph{réseau d'administration hors bande}\index{réseau d'administration hors bande} est un réseau distinct, avec ses propres commutateurs et
liaisons, par lequel une société administre ses appareils et ses serveurs, afin que la configuration, la supervision et l'accès à distance ne
partagent ni un port ni une file avec le trafic de trading, et continuent de fonctionner quand le réseau de trading ne le fait pas. Un \emph{câble à connexion directe}\index{câble à connexion directe}
est un assemblage de câble en cuivre dont les émetteurs-récepteurs sont intégrés à ses extrémités, utilisé pour les liaisons courtes
à l'intérieur d'une baie.
\end{definition}

\begin{method}[Concevoir le réseau d'une cage de trading]\label{met:nw:switches-and-layer-1-devices:design}
\begin{enumerate}
\item Lister les points de remise~: les lignes de données de marché de chaque plateforme (A et B), ses connexions de saisie d'ordres et les liaisons sortantes de la société.
\item Affecter aux données de marché un dispositif de réplication (une \omterm{def:nw:switches-and-layer-1-devices:l1}{duplication de couche 1}, ou un commutateur cut-through lorsque le filtrage par groupe compte), avec
  la ligne A et la ligne B sur des dispositifs distincts, afin qu'une seule panne ne fasse jamais taire les deux.
\item Affecter aux ordres un chemin qui traverse le moins possible de dispositifs~: un \omterm{def:nw:switches-and-layer-1-devices:l1}{multiplexeur de couche 1} ou un commutateur cut-through par point de remise, et
  un deuxième point de remise sur un second dispositif.
\item Placer un tap sur chaque point de remise, dans les deux sens, et acheminer les copies vers les dispositifs de capture~; placer la référence de temps
  (chapitre~4) dans la même cage.
\item Garder l'administration hors bande.
\item Vérifier la conception en retirant tour à tour chaque dispositif et chaque câble~: chaque serveur doit encore recevoir au moins une ligne et
  atteindre au moins un point de remise d'ordres.
\end{enumerate}
\end{method}

\begin{omfigure}
\begin{tikzpicture}[font=\footnotesize,b/.style={draw,rounded corners=2pt,minimum height=0.6cm,align=center,inner sep=2pt},>=Stealth]
\node[b,fill=gray!10] (A) at (0,2.4) {point de remise A};
\node[b,fill=gray!10] (B) at (0,1.2) {point de remise B};
\node[b,fill=gray!10] (O) at (0,-0.4) {point de remise d'ordres 1};
\node[b,fill=gray!10] (P) at (0,-1.6) {point de remise d'ordres 2};
\foreach \n/\y in {tA/2.4,tB/1.2,tO/-0.4,tP/-1.6} {\node[b,fill=omExo!15,minimum width=0.9cm] (\n) at (2.7,\y) {tap};}
\node[b,fill=omDef!20] (la) at (5.2,2.4) {duplication de\\ couche 1};
\node[b,fill=omDef!20] (lb) at (5.2,1.2) {duplication de\\ couche 1};
\node[b,fill=omMeth!20] (m1) at (5.2,-0.4) {multiplexeur};
\node[b,fill=omMeth!20] (m2) at (5.2,-1.6) {multiplexeur};
\node[b,minimum height=1.6cm,minimum width=1.6cm] (s1) at (9.4,1.7) {serveur 1};
\node[b,minimum height=1.6cm,minimum width=1.6cm] (s2) at (9.4,-0.9) {serveur 2};
\node[b,fill=gray!8] (cap) at (2.7,-3.0) {capture et référence de temps};
\foreach \a/\b in {A/tA,B/tB,O/tO,P/tP,tA/la,tB/lb,tO/m1,tP/m2} {\draw[thick] (\a) -- (\b);}
\foreach \a in {la,lb} {\draw[->,omDef,thick] (\a.east) -- (s1.west); \draw[->,omDef,thick] (\a.east) -- (s2.west);}
\foreach \a in {m1,m2} {\draw[<-,omMeth,thick] (\a.east) -- (s1.west); \draw[<-,omMeth,thick] (\a.east) -- (s2.west);}
\draw[gray,dashed] (tP) -- (cap);
\end{tikzpicture}

\omcaption{La conception de couche 1 de la cage du chapitre. Chaque ligne de plateforme passe par un tap, puis par sa propre \omterm{def:nw:switches-and-layer-1-devices:l1}{duplication de couche 1}, qui la copie vers
les deux serveurs~; les ordres de chaque serveur passent par deux multiplexeurs vers deux points de remise d'ordres. Chaque point de remise est muni d'un tap (une seule
copie de tap est représentée)~; l'administration est hors bande et n'est pas dessinée. Retirer un seul dispositif ou un seul câble laisse à chaque serveur une ligne et
un chemin d'ordres.}\label{fig:nw:switches-and-layer-1-devices:cage}
\end{omfigure}

Le \Cref{tab:nw:switches-and-layer-1-devices:designs} compare trois conceptions de la même cage, deux serveurs et les mêmes câbles
(30~mètres d'interconnexion vers chaque point de remise, 3~mètres à l'intérieur de la cage), calculé avec \texttt{firm.cagenet}. La conception sur étagère
place deux commutateurs store-and-forward en série~; la conception cut-through donne à chaque ligne son propre commutateur~; la conception de couche 1 est
la \cref{fig:nw:switches-and-layer-1-devices:cage}.

\begin{table}[ht]
\centering\small
\begin{tabular}{@{}>{\raggedright\arraybackslash}p{3.8cm}>{\centering\arraybackslash}p{2.1cm}>{\centering\arraybackslash}p{2.1cm}>{\centering\arraybackslash}p{2.1cm}>{\raggedright\arraybackslash}p{4.3cm}@{}}
\toprule
Conception & Flux entrant (ns) & Ordre sortant (ns) & Somme (ns) & Points de défaillance uniques \\
\midrule
Commutateurs store-and-forward & 1\,439.8 & 1\,430.2 & 2\,870.0 & les deux commutateurs, trois câbles \\
Commutateurs cut-through & 560.0 & 556.8 & 1\,116.7 & aucun \\
Dispositifs de couche 1 & 262.8 & 294.6 & 557.3 & aucun \\
\bottomrule
\end{tabular}
\caption{Trois conceptions de cage~: latence depuis le premier bit au point de remise de la plateforme jusqu'au dernier bit au serveur pour une trame de
flux de 104~octets, et du serveur au point de remise des ordres pour un ordre de 100~octets, à \qty{10}{\giga\bit\per\second}, \omterm{def:nw:networking-for-trading:egress}{file de sortie} vide. Données~:
\texttt{nw\_cage.design\_table()} sur \texttt{firm.cagenet}.}\label{tab:nw:switches-and-layer-1-devices:designs}
\end{table}

Le chemin du flux de la conception de couche 1 dure 263~nanosecondes, dont 176 pour les 36~mètres de fibre et 83 pour la
sérialisation propre à la trame~: les dispositifs n'ajoutent que 4~nanosecondes. Son chemin d'ordres est plus lent que son chemin de flux de 39~nanosecondes du multiplexeur, moins les
4~nanosecondes de la \omterm{def:nw:switches-and-layer-1-devices:l1}{duplication de couche 1}, et plus lent encore lorsque deux serveurs répondent au même événement. Les nanosecondes suivantes ne sont plus dans les dispositifs, mais dans
les câbles et dans les serveurs.

\section{Tutoriel~: trois cages}\label{tut:nw:switches-and-layer-1-devices:tut}

\begin{tutorial}
\textbf{Objectif.} Modéliser le réseau d'une cage, calculer ses chemins et ses points de défaillance uniques, et faire concourir des ordres dans un multiplexeur.
\textbf{État final~:} \cref{tab:nw:switches-and-layer-1-devices:designs} et \cref{fig:nw:switches-and-layer-1-devices:curve,fig:nw:switches-and-layer-1-devices:mux}.
\begin{tutsteps}
\item \textbf{Dispositifs.} \texttt{nw\_cage.DEVICES} contient les latences des fiches techniques et l'unique valeur de modèle~;
  \texttt{forwarding\_curve()} évalue la \cref{prop:nw:switches-and-layer-1-devices:add} par taille de trame.
\item \textbf{Cages.} \texttt{design(name)} construit chaque cage à partir des dispositifs et des câbles de \texttt{firm.cagenet}~; \texttt{Cage.path} trouve le
  chemin et \texttt{breakdown} décompose sa latence (\cref{lst:nw:switches-and-layer-1-devices:breakdown}).
\item \textbf{Défaillances.} \texttt{Cage.single\_points(feeds, orders, servers)} retire chaque élément à tour de rôle~;
  \texttt{design\_table()} écrit le tableau.
\item \textbf{Contention.} \texttt{mux\_contention(k, jitter)} fait concourir $k$ ordres dans un même multiplexeur sur 20\,000 événements
  (\cref{lst:nw:switches-and-layer-1-devices:mux})~; \texttt{python fig\_cage.py} écrit les données des graphiques.
\end{tutsteps}
\textbf{Pour aller plus loin.} Remplacer les interconnexions par 150~mètres et observer quel terme du tableau domine~; donner à chaque serveur
  sa propre entrée de multiplexeur à partir d'une seconde \omterm{def:nw:networking-for-trading:frame}{carte réseau} et mesurer comment la contention évolue lorsque les quatre stratégies les plus rapides de la société
  partagent deux multiplexeurs au lieu d'un.
\end{tutorial}

\omcode{code/firm/cagenet/firm_cagenet.py}{91}{101}{La latence d'un chemin~: une sérialisation, la fibre et ce qu'ajoute chaque dispositif.}\label{lst:nw:switches-and-layer-1-devices:breakdown}

\omcode{code/networks/02-switches-and-layer-1-devices/python/nw_cage.py}{77}{88}{Ordres de $k$ serveurs passant par un même multiplexeur~:
chacun démarre lorsqu'il est arrivé et que le précédent est terminé.}\label{lst:nw:switches-and-layer-1-devices:mux}

\begin{omfigure}
\begin{tikzpicture}
\begin{axis}[width=11.5cm,height=5.2cm,xmode=log,xmin=20,xmax=2000,ymin=0,ymax=320,
  xlabel={écart-type des délais de réponse des serveurs (ns, échelle logarithmique)},ylabel={attente moyenne d'un ordre (ns)},
  xtick={25,50,100,200,400,800,1600},xticklabels={25,50,100,200,400,800,1\,600},
  tick label style={font=\small},label style={font=\small},grid=major,grid style={gray!25},
  legend columns=3,legend style={font=\footnotesize,at={(0.5,-0.3)},anchor=north,draw=none,fill=none,
  /tikz/every even column/.append style={column sep=8pt}},table/col sep=comma]
\addplot[omDef,very thick,mark=*,mark size=1.4pt] table[x=jitter_ns,y=k2]{figdata/networks/02-switches-and-layer-1-devices/mux.csv};
\addplot[omMeth,very thick,mark=square*,mark size=1.4pt] table[x=jitter_ns,y=k4]{figdata/networks/02-switches-and-layer-1-devices/mux.csv};
\addplot[omThm,very thick,mark=triangle*,mark size=1.8pt] table[x=jitter_ns,y=k8]{figdata/networks/02-switches-and-layer-1-devices/mux.csv};
\legend{2 serveurs,4 serveurs,8 serveurs}
\end{axis}
\end{tikzpicture}

\omcaption{Simulation~: $k$ serveurs envoient chacun un ordre de 100~octets en réponse au même événement, par un même
multiplexeur à \qty{10}{\giga\bit\per\second}~; délai moyen d'attente en file d'un ordre en fonction de la dispersion des temps de réponse des serveurs,
20\,000 événements. Des réponses très groupées font la queue les unes derrière les autres. Données~: \texttt{fig\_cage.py}
(graine fixée).}\label{fig:nw:switches-and-layer-1-devices:mux}
\end{omfigure}

\section{Construction~: le modèle de réseau de cage}\label{bld:nw:switches-and-layer-1-devices:cagenet}

\begin{build}
\bfield{Rôle} La cage de la société sous forme de données~: ce qui est relié à quoi, par quel dispositif, sur quelle longueur de fibre, afin que la latence de chaque chemin
et chaque point de défaillance unique soient calculés plutôt que mémorisés. Le chapitre~5 s'y branche et le chapitre~29 en chiffre le coût.
\bfield{Interface} \texttt{firm\_cagenet}~: \texttt{Device(name, kind, fabric\_ns)} avec les types \texttt{handoff}, \texttt{server},
\texttt{l1}, \texttt{mux}, \texttt{cut-through}, \texttt{store-and-forward}, \texttt{tap}~; \texttt{Cable(a, b, length\_m, n\_g)}~;
\texttt{Cage(devices, cables)} avec \texttt{path}, \texttt{latency\_ns}, \texttt{breakdown},
\texttt{single\_points(feeds, orders, servers)} et \texttt{untapped(handoffs)}.
\bfield{Règles} La latence se mesure du premier bit à la source jusqu'au dernier bit à la destination~; le délai propre à chaque dispositif provient de
\texttt{firm.netsim.forward\_ns}~; les \omterm{def:nw:switches-and-layer-1-devices:tap}{taps réseau} n'ajoutent que de la fibre~; les chemins ne traversent jamais un serveur ni un point de remise~; un type de dispositif inconnu ou un câble vers un
dispositif inconnu est une erreur.
\bfield{Tests d'acceptation} \texttt{code/firm/cagenet/tests/}~: le délai de chaque type de dispositif sur une chaîne d'un seul dispositif~; qu'un chemin ne
traverse pas de serveur~; les points de défaillance uniques et les points de remise sans tap sur une cage construite à la main.
\bfield{Pour aller plus loin} Charge~: acheminer les flux du chapitre~1 à travers la cage avec \texttt{firm.netsim.egress} à chaque sortie et signaler la
pire file d'attente~; un plan de secours listant, pour chaque point de défaillance unique, la reprise de câblage qu'un \omterm{def:nw:switches-and-layer-1-devices:l1}{commutateur de couche 1} effectuerait pour le contourner.
\end{build}

\begin{omsources}
\item Arista, \emph{7130 Connect Series Layer 1 Switch Data Sheet}~; Arista, \emph{7130 MetaMux}, page produit.
\item Cisco, \emph{Nexus 3548-X, 3524-X, 3548-XL and 3524-XL Switches Data Sheet}.
\item Garland Technology, ``Network TAP Split Ratios and Loss Budget'' (2014).
\end{omsources}

\section{Exercices}

\begin{exercise}[$\star$]\label{exo:nw:switches-and-layer-1-devices:1}
Calculer le délai qu'ajoute un commutateur store-and-forward doté d'un temps interne de \qty{300}{\nano\second} à une trame de 512~octets à
\qty{25}{\giga\bit\per\second}. Faire de même pour un commutateur cut-through lisant 64~octets, avec le même temps interne.
\end{exercise}

\begin{exercise}[$\star$]\label{exo:nw:switches-and-layer-1-devices:2}
Un commutateur cut-through reçoit une trame à \qty{10}{\giga\bit\per\second} pour une sortie à \qty{25}{\giga\bit\per\second}, puis une trame à
25 pour une sortie à 10. Déterminer dans quel cas il peut passer en cut-through, et pourquoi.
\end{exercise}

\begin{exercise}[$\star$]\label{exo:nw:switches-and-layer-1-devices:3}
Les deux sens d'une liaison sont recopiés sur un même port. La liaison porte \qty{6}{\giga\bit\per\second} dans chaque sens au moment le plus chargé~;
le port de recopie fonctionne à \qty{10}{\giga\bit\per\second}. Déterminer la part de la copie perdue à ce moment.
\end{exercise}

\begin{exercise}[$\star\star$]\label{exo:nw:switches-and-layer-1-devices:4}
Deux serveurs répondent à un événement avec des délais normaux indépendants d'écart-type \qty{100}{\nano\second}~; leurs ordres de 100~octets passent
par un même multiplexeur à \qty{10}{\giga\bit\per\second}. Déterminer la probabilité que l'ordre le plus tardif doive attendre.
\end{exercise}

\begin{exercise}[$\star\star$]\label{exo:nw:switches-and-layer-1-devices:5}
Dans la conception de couche 1 du \cref{tab:nw:switches-and-layer-1-devices:designs}, décomposer les 262.8~ns du chemin de flux en sérialisation, fibre et équipements.
Lequel raccourcir en premier, et comment~?
\end{exercise}

\begin{exercise}[$\star\star$]\label{exo:nw:switches-and-layer-1-devices:6}
Pourquoi une duplication n'a-t-elle pas besoin de tampon alors qu'un multiplexeur en a besoin~? Qu'est-ce que cela implique pour leur latence et leur gigue~?
\end{exercise}

\begin{exercise}[$\star\star\star$]\label{exo:nw:switches-and-layer-1-devices:7}
\emph{Programmation.} Ajouter un troisième point de remise d'ordres et un troisième multiplexeur à la conception de couche 1, envoyer les ordres du serveur~1 par l'un
et ceux du serveur~2 par les autres, puis recalculer l'attente moyenne de la \cref{fig:nw:switches-and-layer-1-devices:mux} pour quatre serveurs à
\qty{50}{\nano\second} lorsque les quatre se répartissent deux à deux sur deux multiplexeurs.
\end{exercise}

\begin{exercise}[$\star\star\star$]\label{exo:nw:switches-and-layer-1-devices:8}
\emph{Trouver la faille.} «~La ligne A et la ligne B arrivent sur un même \omterm{def:nw:switches-and-layer-1-devices:l1}{commutateur de couche 1}, qui duplique chacune vers nos serveurs. C'est un seul boîtier,
mais les \omterm{def:nw:switches-and-layer-1-devices:l1}{commutateurs de couche 1} n'ont pas de logiciel pouvant planter, donc il n'y a pas de point de défaillance unique.~»
\end{exercise}

\section{Problème~: trois conceptions pour une seule cage}

\begin{problem}[{Problème du week-end --- nanosecondes, pannes et argent}]\label{pb:nw:switches-and-layer-1-devices:1}
Une société de trading choisit entre les trois conceptions du \cref{tab:nw:switches-and-layer-1-devices:designs} pour une cage qui opère sur un signal
issu du flux de la plateforme et répond par un ordre à cette même plateforme. Supposons que la conception cut-through coûte 60\,000 de plus par an
que la conception standard, et la conception de couche 1 150\,000 de plus par an (unités monétaires~; les chiffres sont ceux du problème, non
des cours de marché).

\medskip\noindent\textbf{Partie I --- Les chemins.}
\begin{enumerate}
\item Quel est le chemin de déclenchement de chaque conception, depuis le premier bit de la trame de flux au point de remise jusqu'au dernier bit de l'ordre au
  point de remise d'ordres, en comptant le flux entrant et l'ordre sortant~?
\item Combien la conception cut-through fait-elle économiser par rapport à la conception standard, et la conception de couche 1 par rapport à la conception cut-through~?
\item Dans la conception standard, quelle part du chemin de flux représente la sérialisation des deux commutateurs~?
\item Quelle est la durée, en nanosecondes, des 36~mètres de fibre du chemin de flux de couche 1~?
\end{enumerate}

\medskip\noindent\textbf{Partie II --- Contention.}
\begin{enumerate}[resume]
\item Quatre stratégies de la société partagent le premier multiplexeur de la conception de couche 1 et répondent à un événement à moins de
  \qty{50}{\nano\second} (écart-type) les unes des autres. Quelle est l'attente moyenne d'un ordre (d'après la simulation du chapitre)~?
\item Et celle du dernier des quatre~?
\item La conception reste-t-elle meilleure que la conception cut-through pour le dernier ordre~?
\item Comment répartir les quatre stratégies~?
\end{enumerate}

\medskip\noindent\textbf{Partie III --- Pannes.}
\begin{enumerate}[resume]
\item Énumérer les points de défaillance unique de la conception standard.
\item Pourquoi la conception cut-through n'en a-t-elle aucun alors que chaque ligne ne compte qu'un seul commutateur~?
\item Qu'est-ce qui doit être vrai des points de remise d'ordres de la plateforme pour que la conception de couche 1 n'en ait aucun~?
\item Que devient la latence lorsqu'une des duplications de la conception de couche 1 tombe en panne~?
\end{enumerate}

\medskip\noindent\textbf{Partie IV --- Le verdict.}
\begin{enumerate}[resume]
\item Quel est le coût par nanoseconde et par an du passage de la conception standard à la conception cut-through~?
\item Et du passage de la conception cut-through à la conception de couche 1~?
\item Énoncer le \emph{résultat nommé}~: le chemin de déclenchement de chaque conception et le coût de chaque nanoseconde gagnée.
\item Pour quel type de stratégie le passage à la couche 1 vaut-il son prix~?
\item Qu'une conception de couche 1 abandonne-t-elle de ce qu'un commutateur fournit~?
\item D'où viennent les nanosecondes suivantes, une fois que les équipements n'ajoutent plus que quelques dizaines de nanosecondes~?
\item Quels contrôles de la méthode faut-il appliquer avant de valider la conception~?
\item En une phrase~: quel est l'équipement le plus rapide d'une cage~?
\end{enumerate}
\end{problem}

\section{Questions d'entretien}

\begin{interviewq}[$\star$ \iqroles{développeur}]\label{iq:nw:switches-and-layer-1-devices:1}
Quelle est la différence entre la \omterm{def:nw:switches-and-layer-1-devices:modes}{commutation store-and-forward} et la \omterm{def:nw:switches-and-layer-1-devices:modes}{commutation cut-through}, et quand un commutateur cut-through stocke-t-il puis réexpédie-t-il~?
\end{interviewq}

\begin{interviewq}[$\star$ \iqroles{développeur}]\label{iq:nw:switches-and-layer-1-devices:2}
Qu'est-ce qu'un \omterm{def:nw:switches-and-layer-1-devices:l1}{commutateur de couche 1}, et que ne peut-il pas faire alors qu'un commutateur Ethernet le peut~?
\end{interviewq}

\begin{interviewq}[$\star\star$ \iqroles{développeur}]\label{iq:nw:switches-and-layer-1-devices:3}
Pourquoi les sociétés font-elles passer les ordres par un multiplexeur plutôt que par une duplication, et que se passe-t-il lorsque plusieurs serveurs émettent en même temps~?
\end{interviewq}

\begin{interviewq}[$\star\star$ \iqroles{développeur}]\label{iq:nw:switches-and-layer-1-devices:4}
\omterm{def:nw:switches-and-layer-1-devices:tap}{Tap réseau} ou \omterm{def:nw:switches-and-layer-1-devices:tap}{recopie de port} pour mesurer la latence~? Justifiez votre choix.
\end{interviewq}

\begin{interviewq}[$\star\star\star$ \iqroles{développeur, trader}]\label{iq:nw:switches-and-layer-1-devices:5}
Concevez le réseau d'une cage avec deux lignes vers la plateforme, un point de remise d'ordres et quatre serveurs. Détaillez-moi les cas de panne.
\end{interviewq}

\begin{interviewq}[$\star\star$ \iqroles{développeur}]\label{iq:nw:switches-and-layer-1-devices:6}
Votre commutateur cut-through affiche une latence plus élevée à l'ouverture qu'à midi, pour les mêmes trames. Pourquoi~?
\end{interviewq}
